%! Function Fundamentals — Unit 1, Lesson 1.1 — Group Activity (Answer Key)
\documentclass[10pt]{article}
\usepackage{precalculus-article}
\usepackage{precalculus-key}   % pulls in -boxes; do NOT also load -boxes
\newcommand{\TallMath}[1]{$\displaystyle #1\rule[-1.4em]{0pt}{3.2em}$}

\begin{document}

\pageheader{Unit 1, Lesson 1.1}{Group Activity --- Answer Key}

\vspace{0.12in}

\begin{scenariobox}[The Drone Flight]{plumlight}
\small
A camera drone takes off from a rooftop 2 meters up, climbs, hovers to film, and
then descends to the ground. The function $h$ gives its height in meters
$t$ minutes after launch. Every tier reads this same graph --- no drawing
required.

\begin{center}
\begin{tikzpicture}[x=0.6cm, y=0.5cm]
  \draw[linegray, very thin, step=1] (0,0) grid (8,7);
  \draw[->, thick] (0,0) -- (8.6,0) node[below right] {\small $t$ (min)};
  \draw[->, thick] (0,0) -- (0,7.6) node[above] {\small $h(t)$ (m)};
  \foreach \x in {1,2,3,4,5,6,7,8} \node[below] at (\x,0) {\small \x};
  \foreach \y in {2,4,6} \node[left] at (0,\y) {\small \y};
  \draw[very thick, plum] (0,2) -- (2,6);
  \draw[very thick, plum] (2,6) -- (5,6);
  \draw[very thick, plum] (5,6) -- (8,0);
  \fill[plum] (0,2) circle (2.5pt);
  \fill[plum] (8,0) circle (2.5pt);
\end{tikzpicture}
\end{center}
\end{scenariobox}

\vspace{0.1in}

\boxguard
\begin{tcolorbox}[colback=white, colframe=black!40,
    title=\textbf{Tier R --- Remediate},
    fonttitle=\bfseries, arc=2mm, left=3mm, right=3mm, top=2mm, bottom=2mm]

\begin{enumerate}[itemsep=8pt]
  \item Read outputs from the graph.

        \begin{enumerate}[label=(\alph*), itemsep=4pt]
          \item $h(0) = $ \ans{2} \quad (how high the drone starts)
          \item $h(3) = $ \ans{6}
        \end{enumerate}

  \item The statement $h(1) = 4$ is a full sentence in disguise. Write it in
        words.

        \vspace{0.02in}
        \ansline{One minute after launch, the drone is 4 meters above the ground.}

  \item The drone lands when its height is 0. At what time does it land?
        \quad $t = $ \ans{8} minutes

  \item For this function: \quad Input quantity: \ans{time (minutes)} \quad
        Output quantity: \ans{height of the drone (meters)}
\end{enumerate}
\end{tcolorbox}

\vspace{0.1in}

\boxguard
\begin{tcolorbox}[colback=white, colframe=black!40,
    title=\textbf{Tier A --- Approaching Proficiency},
    fonttitle=\bfseries, arc=2mm, left=3mm, right=3mm, top=2mm, bottom=2mm]

\begin{enumerate}[itemsep=8pt]
  \item Evaluate $h(6)$ from the graph. \quad $h(6) = $ \ans{4}

  \item Solve $h(t) = 4$. \quad $t = $ \ans{1} and $t = $ \ans{6}

        Why is it okay for this equation to have two answers, even though $h$ is
        a function?

        \vspace{0.02in}
        \ansline{Two inputs may share an output (many-to-one is allowed). Each time still has exactly one height, so $h$ is still a function --- the drone just passes 4 meters once going up and once coming down.}

  \item State the domain and range of $h$ as inequalities.

        \vspace{0.05in}
        Domain: \ans{0} $\le t \le$ \ans{8} \quad
        Range: \ans{0} $\le h \le$ \ans{6}

  \item Use the graph to explain how we know $h$ \textbf{is} a function of $t$.

        \vspace{0.02in}
        \ansline{Every time $t$ has exactly one height on the graph --- no input has two outputs.}
\end{enumerate}
\end{tcolorbox}

\vspace{0.1in}

\boxguard
\begin{tcolorbox}[colback=white, colframe=black!40,
    title=\textbf{Tier E --- Extension},
    fonttitle=\bfseries, arc=2mm, left=3mm, right=3mm, top=2mm, bottom=2mm]

\begin{enumerate}[itemsep=8pt]
  \item Flip the machine: is \textbf{time} a function of \textbf{height} for
        this flight? Use the two moments when the drone is 4 meters up in your
        explanation.

        \vspace{0.02in}
        \ansline{No. The single input ``height 4 m'' has two outputs ($t = 1$ and $t = 6$) --- one input with two outputs breaks the function rule.}

  \item A different function is given by the formula $g(x) = \dfrac{12}{x - 5}$.
        Which input must be excluded from the domain, and why?

        \vspace{0.02in}
        \ansline{$x = 5$ must be excluded: it makes the denominator 0, and division by zero is undefined. The domain is all real numbers except 5.}

  \item The formula behind the drone's descent works for any $t$, yet we say
        the domain of $h$ is only $0 \le t \le 8$. Why does the \textit{context}
        shrink the domain?

        \vspace{0.02in}
        \ansline{The flight only exists from launch ($t = 0$) to landing ($t = 8$) --- inputs outside that window describe no real moment of the flight, so the context limits which inputs make sense.}
\end{enumerate}
\end{tcolorbox}

\end{document}
